% ============================================================================
%  ESTADÍSTICA DESCRIPTIVA — Sesión 9
%  Semana 9: 5 al 10 de octubre de 2026
%  2.5 Elaboración de gráficos estadísticos: histograma, polígono y
%      diagrama de caja
%  Duración prevista: 3 horas
% ============================================================================
\documentclass[aspectratio=169,10pt]{beamer}

\usepackage[utf8]{inputenc}
\IfFileExists{spanish.ldf}{\usepackage[spanish,es-noshorthands]{babel}}{}
\usepackage{amsmath,amssymb}
\usepackage{booktabs}
\usepackage{tikz}
\usepackage{listings}
\usepackage{pgfplots}
\pgfplotsset{compat=1.16}

\usetheme{Uniminuto}

\uniasignatura{Estadística descriptiva}
\uniprofesor{Andrés Felipe Puerta Velez}
\unicorreo{andres.puerta-v@uniminuto.edu.co}
\uninrc{CBASBL031 - CBASBL141}
\unisemana{Sesión 9 — Semana 9 · 5 al 10 de octubre de 2026}

\newcommand{\listatiny}{\setbeamerfont{itemize/enumerate body}{size=\tiny}\tiny}

\pgfplotsset{
  uniplot/.style={
    width=8.0cm, height=3.5cm, axis lines=left,
    tick label style={font=\tiny,color=uniGris},
    label style={font=\tiny,color=uniGris},
    title style={font=\scriptsize\bfseries,color=uniAzul},
    every axis plot/.append style={fill=uniAzul!70,draw=uniAzul},
    ymajorgrids, grid style={uniGris!25,very thin},
  }
}

% ---- macro para dibujar una caja horizontal -------------------------------
% \caja{escala x}{min}{Q1}{Me}{Q3}{max}{y}{color}
\newcommand{\cajahor}[8]{%
  \draw[#8,thick] (#2*#1,#7) -- (#3*#1,#7);
  \draw[#8,thick] (#5*#1,#7) -- (#6*#1,#7);
  \draw[#8,thick] (#2*#1,#7-0.16) -- (#2*#1,#7+0.16);
  \draw[#8,thick] (#6*#1,#7-0.16) -- (#6*#1,#7+0.16);
  \draw[#8,thick,fill=#8!18] (#3*#1,#7-0.28) rectangle (#5*#1,#7+0.28);
  \draw[#8,very thick] (#4*#1,#7-0.28) -- (#4*#1,#7+0.28);
}

\newcommand{\fx}[1]{{\ttfamily\color{uniAzul}#1}}

% ------------------------------------------------------------- CÓDIGO R ----
\lstdefinestyle{Rcode}{
  language=R,
  basicstyle=\ttfamily\scriptsize,
  keywordstyle=\color{uniAzul}\bfseries,
  commentstyle=\color{uniGris}\itshape,
  stringstyle=\color{uniTitulo},
  showstringspaces=false, numbers=none, breaklines=true,
  columns=fullflexible,
  backgroundcolor=\color{uniCremaClaro},
  frame=leftline, framerule=1.2pt, rulecolor=\color{uniNaranja},
  xleftmargin=0.35cm, framexleftmargin=0.35cm,
  framexrightmargin=0.15cm, framextopmargin=2pt, framexbottommargin=2pt,
  aboveskip=0.15cm, belowskip=0.15cm
}
\lstdefinestyle{Rmini}{
  language=R,
  basicstyle=\ttfamily\tiny,
  keywordstyle=\color{uniAzul}\bfseries,
  commentstyle=\color{uniGris}\itshape,
  stringstyle=\color{uniTitulo},
  showstringspaces=false, numbers=none, breaklines=true,
  columns=fullflexible,
  backgroundcolor=\color{uniCremaClaro},
  frame=leftline, framerule=1.2pt, rulecolor=\color{uniNaranja},
  xleftmargin=0.35cm, framexleftmargin=0.35cm,
  framexrightmargin=0.15cm, framextopmargin=2pt, framexbottommargin=2pt,
  aboveskip=0.12cm, belowskip=0.12cm
}
\lstdefinestyle{Rout}{
  basicstyle=\ttfamily\tiny\color{black},
  showstringspaces=false, numbers=none, breaklines=true,
  columns=fullflexible,
  backgroundcolor=\color{uniGris!10},
  frame=leftline, framerule=1.2pt, rulecolor=\color{uniGris},
  xleftmargin=0.35cm, framexleftmargin=0.35cm,
  framexrightmargin=0.15cm, framextopmargin=2pt, framexbottommargin=2pt,
  aboveskip=0.12cm, belowskip=0.12cm
}

\begin{document}

\portada

\begin{frame}{Orden del día}
  \begin{enumerate}
    \item Llamado a lista y revisión de los ejercicios propuestos
    \item Repaso: histograma y polígono con las medidas encima
    \item \azul{El resumen de cinco números}
    \item \azul{Diagrama de caja: construcción, bigotes y valores atípicos}
    \item Comparación de grupos con cajas paralelas
    \item Taller guiado y cierre
  \end{enumerate}

  \vspace{0.25cm}
  {\footnotesize\color{uniGris} Sesión de 3 horas. Hoy se dibuja mucho: traiga
  regla y papel cuadriculado.}
\end{frame}

\begin{frame}{Repaso: el histograma con las medidas encima}
  \centering
  \begin{tikzpicture}
    \begin{axis}[uniplot, width=8.6cm, height=3.6cm,
      title={Gasto mensual — Almacén Bello, con $Mo$, $Me$ y $\bar{x}$},
      ybar interval, bar width=1,
      xlabel={Gasto (miles \$)}, ylabel={$f_i$},
      xtick={112,152,192,232,272,312,352}, ymin=0, ymax=13]
      \addplot coordinates {(112,4) (152,7) (192,11) (232,9) (272,6) (312,3)
                            (352,0)};
      \draw[uniTitulo,very thick] (axis cs:218.67,0) -- (axis cs:218.67,12);
      \draw[uniNaranja,very thick,dashed] (axis cs:224.73,0)
                                          -- (axis cs:224.73,12);
      \draw[uniGris,very thick,dotted] (axis cs:227,0) -- (axis cs:227,12);
    \end{axis}
  \end{tikzpicture}

  \vspace{-0.1cm}
  \begin{alertblock}{Las tres líneas}
    \scriptsize
    \destacar{Moda 218,67} (línea continua), \destacar{mediana 224,73}
    (discontinua) y \destacar{media 227,0} (punteada). Están casi pegadas y en
    ese orden: asimetría positiva leve, como calculamos en la sesión 8.
  \end{alertblock}
\end{frame}

\begin{frame}{El resumen de cinco números}
  \footnotesize
  \begin{center}
  \begin{tabular}{@{}lrl@{}}
    \toprule
    \textbf{Valor} & \textbf{Almacén Bello} & \textbf{Qué indica} \\
    \midrule
    Mínimo   & 112 & El cliente que menos gasta \\
    $Q_1$    & 185{,}50 & El 25\% gasta menos que esto \\
    Mediana  & 227{,}50 & La mitad gasta menos que esto \\
    $Q_3$    & 267{,}75 & El 75\% gasta menos que esto \\
    Máximo   & 352 & El cliente que más gasta \\
    \bottomrule
  \end{tabular}
  \end{center}

  \vspace{0.15cm}
  \begin{alertblock}{Por qué cinco números y no uno}
    \scriptsize
    Con estos cinco valores se reconstruye \destacar{el centro, la dispersión y
    la asimetría} de la distribución sin necesidad de ver los 40 datos. El
    diagrama de caja no es más que \azul{el dibujo de esta tabla}.
  \end{alertblock}

  \vspace{0.1cm}
  {\scriptsize\color{uniGris} Nota: aquí usamos los cuartiles de los
  \azul{datos originales}. Si solo se dispone de la tabla agrupada, se usan los
  cuartiles interpolados de la sesión 7.}
\end{frame}

\seccion{Diagrama de caja}

\begin{frame}{Anatomía del diagrama de caja}
  \centering
  \begin{tikzpicture}[scale=1.0]
    \cajahor{0.028}{112}{185.5}{227.5}{267.75}{352}{0}{uniAzul}
    \draw[uniGris,->] (2.8,-0.75) -- (11.0,-0.75);
    \foreach \v in {112,150,200,250,300,352}
      \draw[uniGris] (\v*0.028,-0.7) -- (\v*0.028,-0.8)
        node[below,font=\tiny,color=uniGris] {\v};
    \node[font=\tiny,color=uniTitulo,above] at (112*0.028,0.4) {mínimo};
    \node[font=\tiny,color=uniTitulo,above] at (185.5*0.028,0.4) {$Q_1$};
    \node[font=\tiny,color=uniTitulo,above] at (227.5*0.028,0.62) {$Me$};
    \node[font=\tiny,color=uniTitulo,above] at (267.75*0.028,0.4) {$Q_3$};
    \node[font=\tiny,color=uniTitulo,above] at (352*0.028,0.4) {máximo};
    \draw[uniNaranja,<->] (185.5*0.028,-0.45) -- (267.75*0.028,-0.45)
      node[midway,below,font=\tiny,color=uniNaranja] {$RIC = 82{,}25$};
  \end{tikzpicture}

  \vspace{0.35cm}
  \begin{alertblock}{Qué representa cada parte}
    \scriptsize
    La \destacar{caja} contiene al 50\% central de los datos; su ancho es el
    $RIC$. La \destacar{línea interior} es la mediana. Cada \destacar{bigote}
    cubre aproximadamente el 25\% de los datos. Si un bigote es mucho más
    largo que el otro, hay \azul{asimetría}.
  \end{alertblock}
\end{frame}

\begin{frame}{Construcción paso a paso}
  \footnotesize
  \begin{block}{Los seis pasos}
    \tiny
    \begin{enumerate}\setlength{\itemsep}{-0.03cm}
      \item Ordenar los datos y calcular $Q_1$, $Me$ y $Q_3$.
      \item Calcular $RIC = Q_3 - Q_1$.
      \item Calcular las \azul{vallas}: \ inferior $= Q_1 - 1{,}5\,RIC$; \
            superior $= Q_3 + 1{,}5\,RIC$.
      \item Marcar como \destacar{atípicos} los datos fuera de las vallas.
      \item Dibujar los bigotes hasta el dato \azul{más extremo que aún esté
            dentro} de las vallas —no hasta la valla.
      \item Dibujar los atípicos como puntos sueltos.
    \end{enumerate}
  \end{block}

  \vspace{0.05cm}
  \begin{alertblock}{El error número uno}
    \tiny
    Los bigotes \destacar{no llegan hasta las vallas}: llegan hasta el último
    dato real que cabe dentro de ellas.
  \end{alertblock}
\end{frame}

\begin{frame}{Aplicación al Almacén Bello}
  \scriptsize
  \begin{columns}[T,onlytextwidth]
    \begin{column}{0.46\textwidth}
      \begin{exampleblock}{Los cálculos}
        \scriptsize
        \[ RIC = 267{,}75-185{,}50 = 82{,}25 \]
        \[ 1{,}5 \times 82{,}25 = 123{,}375 \]
        \[ \text{V. inf.} = 185{,}50-123{,}38 = 62{,}12 \]
        \[ \text{V. sup.} = 267{,}75+123{,}38 = 391{,}13 \]
      \end{exampleblock}
    \end{column}
    \begin{column}{0.50\textwidth}
      \begin{block}{Conclusión}
        \scriptsize
        El mínimo (112) y el máximo (352) están \azul{dentro} de las vallas:
        \destacar{no hay valores atípicos}.

        \vspace{0.1cm}
        Los bigotes llegan entonces hasta 112 y hasta 352.
      \end{block}
    \end{column}
  \end{columns}

  \vspace{0.1cm}
  \begin{alertblock}{Leyendo la asimetría en la caja}
    \scriptsize
    Distancia $Q_1$ a $Me$: $227{,}5-185{,}5 = 42{,}0$. \quad
    Distancia $Me$ a $Q_3$: $267{,}75-227{,}5 = 40{,}25$. \quad
    Casi iguales: la mitad central es \azul{prácticamente simétrica}. El bigote
    derecho (84,25) sí es más largo que el izquierdo (73,5): la
    \destacar{cola} derecha es la que produce la asimetría positiva.
  \end{alertblock}
\end{frame}

\begin{frame}{Un caso con valores atípicos}
  \footnotesize
  Tiempo de respuesta de soporte técnico, en minutos ($n=15$):

  \vspace{0.1cm}
  \begin{beamercolorbox}[wd=\textwidth,sep=0.1cm]{block body}
    \centering\scriptsize\ttfamily
    5 \ 7 \ 8 \ 9 \ 10 \ 11 \ 12 \ 13 \ 14 \ 15 \ 16 \ 18 \ 20 \ 45 \ 60
  \end{beamercolorbox}

  \vspace{0.15cm}
  \centering
  \begin{tikzpicture}[scale=1.0]
    \cajahor{0.16}{5}{9}{13}{18}{20}{0}{uniAzul}
    \fill[uniTitulo] (45*0.16,0) circle (2.2pt);
    \fill[uniTitulo] (60*0.16,0) circle (2.2pt);
    \draw[uniGris,->] (0.4,-0.65) -- (10.4,-0.65);
    \foreach \v in {0,10,20,30,40,50,60}
      \draw[uniGris] (\v*0.16,-0.6) -- (\v*0.16,-0.7)
        node[below,font=\tiny,color=uniGris] {\v};
    \draw[uniNaranja,dashed] (31.5*0.16,-0.45) -- (31.5*0.16,0.55)
      node[above,font=\tiny,color=uniNaranja] {valla sup. $=31{,}5$};
  \end{tikzpicture}

  \vspace{0.15cm}
  \begin{alertblock}{Los cálculos}
    \scriptsize
    $Q_1 = 9$, \ $Me = 13$, \ $Q_3 = 18$, \ $RIC = 9$, \ valla superior
    $= 18 + 13{,}5 = 31{,}5$. Los valores \destacar{45 y 60} la superan: son
    \azul{atípicos}. El bigote derecho llega hasta \destacar{20}, el último
    dato admitido.
  \end{alertblock}
\end{frame}

\begin{frame}{Qué hacer con un valor atípico}
  \scriptsize
  \begin{center}
  \begin{tabular}{@{}p{4.2cm}p{7.6cm}@{}}
    \toprule
    \textbf{Posible causa} & \textbf{Qué se hace} \\
    \midrule
    Error de digitación o de unidad &
      Se corrige o se elimina, dejando constancia en el informe \\
    \addlinespace
    Caso legítimo pero excepcional &
      Se conserva y se analiza aparte: suele ser lo más interesante \\
    \addlinespace
    Mezcla de dos poblaciones &
      Se separan los grupos y se describe cada uno \\
    \bottomrule
  \end{tabular}
  \end{center}

  \vspace{0.15cm}
  \begin{alertblock}{La regla profesional}
    \scriptsize
    \destacar{Nunca se borra un dato solo porque estorbe.} En el ejemplo de
    soporte, los casos de 45 y 60 minutos podrían ser justamente los que
    generan las quejas: eliminarlos ocultaría el problema que se quiere
    resolver.
  \end{alertblock}

  \vspace{0.1cm}
  {\tiny\color{uniGris} El criterio $1{,}5 \times RIC$ es una convención (Tukey),
  no una ley física. Con $3 \times RIC$ se marcan los atípicos ``extremos''.}
\end{frame}

\begin{frame}{Comparar grupos con cajas paralelas}
  \centering
  \begin{tikzpicture}[scale=1.0]
    \cajahor{0.32}{4}{7}{10}{13}{18}{0.9}{uniAzul}
    \cajahor{0.32}{2}{5}{9}{17}{28}{0}{uniNaranja}
    \node[font=\tiny,color=uniAzul,left] at (0.3,0.9) {Bello};
    \node[font=\tiny,color=uniNaranja,left] at (0.3,0) {Itagüí};
    \draw[uniGris,->] (0.4,-0.6) -- (9.9,-0.6);
    \foreach \v in {0,5,10,15,20,25,30}
      \draw[uniGris] (\v*0.32,-0.55) -- (\v*0.32,-0.65)
        node[below,font=\tiny,color=uniGris] {\v};
    \node[font=\tiny,color=uniGris] at (5.0,-1.05) {Tiempo de espera (minutos)};
  \end{tikzpicture}

  \vspace{0.35cm}
  \begin{alertblock}{Lo que se lee de un vistazo}
    \scriptsize
    Las medianas son parecidas (10 y 9 minutos), pero la caja de
    \destacar{Itagüí es mucho más ancha}: $RIC = 12$ frente a 6. El servicio
    de Bello es \azul{más predecible}. Además la caja de Itagüí es asimétrica
    hacia arriba: hay clientes que esperan mucho más de lo habitual.
  \end{alertblock}

  \vspace{0.05cm}
  {\tiny\color{uniGris} Este es el gráfico que se lleva a una junta cuando hay
  que comparar sucursales, turnos, meses o productos.}
\end{frame}

\begin{frame}{Histograma o caja: ¿cuál usar?}
  \scriptsize
  \begin{center}
  \begin{tabular}{@{}p{2.4cm}p{4.6cm}p{4.6cm}@{}}
    \toprule
    & \textbf{Histograma} & \textbf{Diagrama de caja} \\
    \midrule
    Muestra & La forma completa: picos, huecos, bimodalidad &
      Cinco números: centro, dispersión y atípicos \\
    \addlinespace
    Oculta & Los valores atípicos individuales &
      La forma: dos distribuciones muy distintas pueden dar cajas iguales \\
    \addlinespace
    Comparar grupos & Difícil: hay que superponer o repetir &
      Muy fácil: cajas paralelas \\
    \addlinespace
    Necesita & Bastantes datos y clases bien elegidas &
      Solo los cuartiles; funciona con pocos datos \\
    \bottomrule
  \end{tabular}
  \end{center}

  \vspace{0.1cm}
  \begin{alertblock}{En la práctica}
    \scriptsize
    \destacar{Se usan juntos.} El histograma para entender la variable; la
    caja para comparar grupos y detectar atípicos. En un informe, primero el
    histograma de la variable y después las cajas por segmento.
  \end{alertblock}
\end{frame}

\begin{frame}{Ejercicio 1 — en clase}
  \begin{alertblock}{En parejas (20 minutos)}
    \footnotesize
    Monto de 12 facturas de un restaurante (miles de pesos):

    \vspace{0.15cm}
    \begin{beamercolorbox}[wd=0.9\textwidth,sep=0.1cm]{block body}
      \centering\scriptsize\ttfamily
      21 \ 33 \ 14 \ 25 \ 19 \ 48 \ 22 \ 30 \ 16 \ 27 \ 18 \ 24
    \end{beamercolorbox}

    \vspace{0.2cm}
    \begin{enumerate}\setlength{\itemsep}{-0.05cm}
      \item Ordene los datos y obtenga el resumen de cinco números.
      \item Calcule el $RIC$ y las dos vallas.
      \item ¿Hay valores atípicos? ¿Cuáles?
      \item ¿Hasta dónde llegan los bigotes?
      \item Dibuje el diagrama de caja a escala.
      \item ¿Qué asimetría muestra? Justifique con las distancias.
    \end{enumerate}
  \end{alertblock}
\end{frame}

\begin{frame}{Ejercicio 1 — solución}
  \scriptsize
  Ordenados: \ttfamily 14 \ 16 \ 18 \ 19 \ 21 \ 22 \ 24 \ 25 \ 27 \ 30 \ 33
  \ 48 \rmfamily \quad ($n = 12$)

  \vspace{0.1cm}
  \begin{columns}[T,onlytextwidth]
    \begin{column}{0.47\textwidth}
      \begin{exampleblock}{Puntos 1 a 4}
        \scriptsize
        pos$(Q_1) = \frac{13}{4} = 3{,}25 \Rightarrow
        Q_1 = 18+0{,}25(19-18) = \destacar{18{,}25}$

        pos$(Me) = 6{,}5 \Rightarrow Me = \frac{22+24}{2} = \destacar{23}$

        pos$(Q_3) = 9{,}75 \Rightarrow Q_3 = 27+0{,}75(30-27) = \destacar{29{,}25}$
        \[ RIC = 11{,}0 \qquad 1{,}5\,RIC = 16{,}5 \]
        Vallas: $1{,}75$ y $\destacar{45{,}75}$ \\
        Atípico: \destacar{48}. \ Bigotes: de 14 a 33.
      \end{exampleblock}
    \end{column}
    \begin{column}{0.49\textwidth}
      \centering
      \begin{tikzpicture}[scale=0.95]
        \cajahor{0.115}{14}{18.25}{23}{29.25}{33}{0}{uniAzul}
        \fill[uniTitulo] (48*0.115,0) circle (2pt);
        \draw[uniGris,->] (0.8,-0.6) -- (6.0,-0.6);
        \foreach \v in {10,20,30,40,50}
          \draw[uniGris] (\v*0.115,-0.55) -- (\v*0.115,-0.65)
            node[below,font=\tiny,color=uniGris] {\v};
      \end{tikzpicture}
    \end{column}
  \end{columns}

  \vspace{0.1cm}
  \begin{alertblock}{Punto 6}
    \tiny
    $Me-Q_1 = 4{,}75$ y $Q_3-Me = 6{,}25$: la mitad superior de la caja es más
    ancha. Sumado al atípico de 48, la distribución es
    \destacar{asimétrica positiva}.
  \end{alertblock}
\end{frame}

\begin{frame}{Ejercicio 2 — en clase (taller principal)}
  \begin{alertblock}{En grupos de 3 (25 minutos)}
    \scriptsize
    Dos turnos de una misma bodega registran las unidades despachadas por hora:

    \vspace{0.15cm}
    \begin{beamercolorbox}[wd=0.94\textwidth,sep=0.1cm]{block body}
      \centering\scriptsize\ttfamily
      Turno mañana: \ 30 \ 32 \ 35 \ 36 \ 38 \ 40 \ 41 \ 43 \ 45 \ 60 \\
      Turno tarde: \ \ 22 \ 28 \ 33 \ 36 \ 38 \ 39 \ 42 \ 48 \ 52 \ 55
    \end{beamercolorbox}

    \vspace{0.2cm}
    \begin{enumerate}\setlength{\itemsep}{-0.07cm}
      \item Obtenga el resumen de cinco números de cada turno.
      \item Calcule $RIC$ y vallas de cada uno; identifique atípicos.
      \item Dibuje las dos cajas en paralelo, con la misma escala.
      \item ¿Qué turno despacha más, en mediana?
      \item ¿Qué turno es más consistente? Justifique con el $RIC$.
      \item Si usted fuera el jefe de bodega, ¿qué decisión tomaría y con qué
            argumento?
    \end{enumerate}
  \end{alertblock}
\end{frame}

\begin{frame}{Ejercicio 2 — solución}
  \scriptsize
  \begin{columns}[T,onlytextwidth]
    \begin{column}{0.44\textwidth}
      \centering
      \begin{tabular}{@{}lrr@{}}
        \toprule
        & \textbf{Mañana} & \textbf{Tarde} \\
        \midrule
        Mínimo & 30 & 22 \\
        $Q_1$  & 34{,}25 & 31{,}75 \\
        $Me$   & 39{,}0  & 38{,}5 \\
        $Q_3$  & 43{,}5  & 49{,}0 \\
        Máximo & 60 & 55 \\
        \midrule
        $RIC$  & 9{,}25 & 17{,}25 \\
        V. sup.& 57{,}38 & 74{,}88 \\
        Atípico & \destacar{60} & --- \\
        \bottomrule
      \end{tabular}
    \end{column}
    \begin{column}{0.53\textwidth}
      \centering
      \begin{tikzpicture}[scale=0.92]
        \cajahor{0.105}{30}{34.25}{39}{43.5}{45}{0.85}{uniAzul}
        \fill[uniTitulo] (60*0.105,0.85) circle (2pt);
        \cajahor{0.105}{22}{31.75}{38.5}{49}{55}{0}{uniNaranja}
        \node[font=\tiny,color=uniAzul,left] at (2.0,0.85) {Mañana};
        \node[font=\tiny,color=uniNaranja,left] at (2.0,0) {Tarde};
        \draw[uniGris,->] (2.05,-0.55) -- (6.7,-0.55);
        \foreach \v in {20,30,40,50,60}
          \draw[uniGris] (\v*0.105,-0.5) -- (\v*0.105,-0.6)
            node[below,font=\tiny,color=uniGris] {\v};
      \end{tikzpicture}
    \end{column}
  \end{columns}

  \vspace{0.05cm}
  \begin{alertblock}{Puntos 4 a 6}
    \tiny
    Las medianas son casi iguales (39 y 38,5): \destacar{ninguno despacha
    claramente más}. Pero el turno de la \azul{mañana} tiene un $RIC$ de 9,25
    frente a 17,25: es casi el doble de consistente. Decisión razonable: revisar
    qué causa la variabilidad de la tarde (¿personal cambiante? ¿picos de
    pedidos?) e investigar la hora de 60 unidades de la mañana, que fue
    atípica.
  \end{alertblock}
\end{frame}

\begin{frame}{Buenas prácticas al presentar gráficos}
  \scriptsize
  \begin{center}
  \begin{tabular}{@{}p{4.4cm}p{7.4cm}@{}}
    \toprule
    \textbf{Práctica} & \textbf{Por qué} \\
    \midrule
    Título que diga qué, quién y cuándo &
      ``Gasto mensual de 40 clientes, Almacén Bello, agosto 2026'' \\
    Ejes rotulados con unidades &
      Sin unidad, el número no significa nada \\
    Misma escala al comparar &
      Dos cajas en escalas distintas no son comparables \\
    Indicar $n$ en cada grupo &
      Una caja con 5 datos no vale lo mismo que una con 500 \\
    Señalar los atípicos, no esconderlos &
      Suelen ser la información más valiosa \\
    Un mensaje por gráfico &
      Si el gráfico necesita tres párrafos de explicación, sobra \\
    \bottomrule
  \end{tabular}
  \end{center}
\end{frame}

\begin{frame}[fragile]{El diagrama de caja, en R}
  \scriptsize
\begin{lstlisting}[style=Rmini]
boxplot(gasto, horizontal = TRUE, col = "steelblue",
        main = "Gasto mensual de 40 clientes", xlab = "Miles de pesos")

fivenum(gasto)              # los cinco numeros que dibuja el grafico
boxplot.stats(gasto)$stats  # bigote inf, Q1, mediana, Q3, bigote sup
boxplot.stats(gasto)$out    # valores atipicos

# Cajas paralelas: la formula se lee "gasto SEGUN sucursal"
boxplot(gasto ~ sucursal, col = "steelblue",
        main = "Gasto por sucursal", xlab = "Sucursal", ylab = "Miles de pesos")
\end{lstlisting}
\begin{lstlisting}[style=Rout]
[1] 112.0 187.0 227.5 266.5 352.0
[1] 112.0 187.0 227.5 266.5 352.0
numeric(0)
\end{lstlisting}

  \vspace{-0.1cm}
  \begin{alertblock}{Dos detalles finos}
    \tiny
    \fx{numeric(0)} significa \destacar{ningún atípico}, igual que concluimos a
    mano. Y ojo: \fx{fivenum()} usa las \azul{bisagras de Tukey} (187 y 266,5)
    mientras \fx{quantile()} usa otra convención (188,5 y 265,25). El dibujo se
    hace con las primeras.
  \end{alertblock}
\end{frame}

\begin{frame}{Cierre de la sesión}
  \begin{block}{Lo que debe quedar claro hoy}
    \scriptsize
    \begin{itemize}\setlength{\itemsep}{0.02cm}
      \item El \azul{resumen de cinco números} describe una distribución sin
            mostrar los datos.
      \item Caja $=$ 50\% central; línea $=$ mediana; bigotes $\approx$
            25\% a cada lado.
      \item Vallas $= Q_1 - 1{,}5\,RIC$ y $Q_3 + 1{,}5\,RIC$; los bigotes
            llegan al \destacar{último dato dentro} de ellas.
      \item Un atípico se investiga, no se borra.
      \item Cajas paralelas: la herramienta para \azul{comparar grupos}.
    \end{itemize}
  \end{block}

  \vspace{0.1cm}
  \begin{alertblock}{Trabajo independiente}
    \scriptsize
    Resolver los 20 ejercicios de la lámina siguiente. En la sesión 10 pasamos
    de calcular a \destacar{interpretar}: patrones, tendencias y redacción del
    informe estadístico.
  \end{alertblock}
\end{frame}

\begin{frame}{Ejercicios propuestos}
  \vspace{-0.1cm}
  \begin{beamercolorbox}[wd=\textwidth,sep=0.10cm]{block body}
    \centering\scriptsize
    \textbf{1--7:} cinco números y vallas. \quad \textbf{8--14:} construcción.
    \quad \textbf{15--20:} lectura e interpretación.
  \end{beamercolorbox}
  \listatiny
  \begin{columns}[T,onlytextwidth]
    \begin{column}{0.47\textwidth}
      \begin{enumerate}\setlength{\itemsep}{0.16cm}
        \item Cinco números de: 3, 5, 6, 8, 9, 11, 14, 20
        \item $RIC$ de ese conjunto
        \item Vallas de ese conjunto
        \item ¿Hay atípicos? ¿Cuáles?
        \item Con $Q_1=40$ y $Q_3=70$: $RIC$ y vallas
        \item ¿Es atípico un dato de 120 en ese caso?
        \item ¿Y uno de 10?
        \item ¿Hasta dónde llega el bigote superior si el mayor dato dentro de
              la valla es 108?
        \item Dibuje la caja de: 2, 4, 5, 5, 6, 7, 9
        \item ¿Qué porcentaje de datos hay dentro de la caja?
      \end{enumerate}
    \end{column}
    \begin{column}{0.50\textwidth}
      \begin{enumerate}\setlength{\itemsep}{0.16cm}
        \setcounter{enumi}{10}
        \item ¿Qué porcentaje cubre cada bigote?
        \item Si la mediana está pegada a $Q_1$, ¿qué asimetría hay?
        \item ¿Qué caja es más dispersa: $RIC=5$ o $RIC=15$?
        \item ¿Puede un diagrama de caja no tener bigote de un lado?
        \item Dos cajas con igual mediana y distinto $RIC$: ¿qué concluye?
        \item ¿Por qué no se borra un atípico automáticamente?
        \item ¿Qué muestra el histograma que no muestra la caja?
        \item ¿Qué muestra la caja que no muestra el histograma?
        \item En el Almacén Bello, ¿por qué no hay atípicos?
        \item Escriba una frase de interpretación de la caja del turno tarde
      \end{enumerate}
    \end{column}
  \end{columns}
\end{frame}

\begin{frame}{Respuestas}{Autoevaluación de los ejercicios propuestos}
  \listatiny
  \begin{columns}[T,onlytextwidth]
    \begin{column}{0.47\textwidth}
      \begin{enumerate}\setlength{\itemsep}{0.16cm}
        \item Mín 3; $Q_1=5{,}25$; $Me=8{,}5$; $Q_3=13{,}25$; máx 20
        \item $RIC = 8{,}0$
        \item $-6{,}75$ y $25{,}25$
        \item No hay atípicos
        \item $RIC = 30$; vallas $-5$ y $115$
        \item Sí: $120 > 115$
        \item No: $10 > -5$
        \item Hasta 108
        \item $Q_1=4$; $Me=5$; $Q_3=7$; mín 2; máx 9
        \item El 50\%
      \end{enumerate}
    \end{column}
    \begin{column}{0.50\textwidth}
      \begin{enumerate}\setlength{\itemsep}{0.16cm}
        \setcounter{enumi}{10}
        \item Aproximadamente el 25\% cada uno
        \item Positiva: la cola larga está arriba
        \item La de $RIC = 15$
        \item Sí: si $Q_1$ coincide con el mínimo, por ejemplo
        \item Mismo centro, distinta consistencia
        \item Puede ser un caso legítimo e informativo
        \item La forma: picos, huecos, bimodalidad
        \item Los cuartiles y los atípicos, y permite comparar grupos
        \item Porque 112 y 352 caen dentro de $[62{,}12;\ 391{,}13]$
        \item Respuesta abierta: mediana 38,5 con alta variabilidad
      \end{enumerate}
    \end{column}
  \end{columns}
\end{frame}

\begin{frame}{Referencias}
  \footnotesize
  \begin{itemize}
    \item[{[1]}] Monroy Saldívar, S. (2008). \textit{Estadística descriptiva}.
      Instituto Politécnico Nacional.
      \texttt{https://elibro.net/es/lc/uniminuto/titulos/74722}
    \item[{[2]}] Ordóñez Fernández, F. F. \& González Fernández, J. (2021).
      \textit{Estadística descriptiva paso a paso}. Ediciones Pirámide.
      \texttt{https://elibro.net/es/lc/uniminuto/titulos/215449}
    \item[{[3]}] Llinás Solano, H. (2014). \textit{Introducción a la
      estadística con aplicaciones en ciencias sociales}. Universidad del
      Norte. \texttt{https://elibro.net/es/lc/uniminuto/titulos/70062}
  \end{itemize}

  \vspace{0.35cm}
  \begin{center}
    {\Large\color{uniAzul}\textbf{¿Preguntas?}}
  \end{center}
\end{frame}

\end{document}
